\section{Обеспечение надёжности и безопасности сервиса}
\thispagestyle{plain}

В настоящей главе положения теории надёжности применяются
к собственной архитектуре сервиса EduQuiz. Выполняется расчёт
показателей надёжности, обосновываются стратегии резервирования
и реакции на отказы, описываются мероприятия по обеспечению
информационной безопасности и приводятся результаты нагрузочного
тестирования. Тем самым достигается замкнутая инженерно-педагогическая
петля настоящей работы: дисциплина <<Надёжность>>, преподаваемая
на кафедре АСУ, применена к сервису, который сам выступает
средством её преподавания.

\subsection{Структурная схема надёжности и расчёт показателей}

Архитектура сервиса декомпозируется на следующие узлы:
$N_1$ --- сервер приложения (Go API),
$N_2$ --- база данных (PostgreSQL),
$N_3$ --- кеш и шина событий (Redis),
$N_4$ --- обратный прокси (Caddy с TLS-терминацией),
$N_5$ --- виртуальная машина хостинг-провайдера
(физический и сетевой уровень).

Для каждого узла определены интенсивность отказов $\lambda_i$
и среднее время восстановления $T_{\text{в}i}$. Значения принимаются
по эксплуатационным данным аналогичных систем и SLA провайдера
(табл.~\ref{tab:reliability-nodes}). Для произвольного узла
вероятность безотказной работы за время $t$ при экспоненциальном
законе:
\[
    P_i(t) = e^{-\lambda_i t}.
\]
Средняя наработка до отказа $\text{MTBF}_i = 1 / \lambda_i$.
Коэффициент готовности
\[
    k_{\text{г}\, i} = \frac{\text{MTBF}_i}{\text{MTBF}_i + T_{\text{в}\, i}}.
\]

\begin{table}[!htb]
\fontsize{11pt}{13pt}\selectfont
\centering
\caption{Показатели надёжности узлов системы (расчётные)}
\label{tab:reliability-nodes}
\smallskip
\begin{tabularx}{\linewidth}{|l|c|c|c|c|}
    \hline
    \textbf{Узел} & $\lambda$, ч$^{-1}$ & MTBF, ч & $T_{\text{в}}$, ч & $k_{\text{г}}$ \\
    \hline
        $N_1$ Go API           & $5 \cdot 10^{-5}$ & $20{,}000$ & 0{,}2 & 0{,}99999 \\
    \hline
        $N_2$ PostgreSQL       & $2 \cdot 10^{-5}$ & $50{,}000$ & 0{,}5 & 0{,}99999 \\
    \hline
        $N_3$ Redis            & $5 \cdot 10^{-5}$ & $20{,}000$ & 0{,}2 & 0{,}99999 \\
    \hline
        $N_4$ Caddy            & $5 \cdot 10^{-5}$ & $20{,}000$ & 0{,}1 & 0{,}99999 \\
    \hline
        $N_5$ VM (Selectel)    & $1 \cdot 10^{-4}$ & $10{,}000$ & 1{,}0 & 0{,}9999 \\
    \hline
\end{tabularx}
\end{table}

С точки зрения структурной схемы надёжности узлы $N_1$--$N_4$
объединены последовательно (отказ любого приводит к отказу системы),
$N_5$ выступает общей платформой. Системный коэффициент готовности
без резервирования:
\[
    k_{\text{г}}^{\text{сист}} = \prod_{i=1}^{5} k_{\text{г}\, i} \approx 0{,}9998.
\]
Без специальных мер ожидаемое время недоступности
${\sim} 1 - k_{\text{г}}^{\text{сист}} \approx 1{,}75$ часа в год,
что превышает целевой SLA $0{,}999$ ($\sim$ 8{,}77 часов в год).

\subsection{Резервирование и стратегия отказоустойчивости}

Для повышения коэффициента готовности применяется горячее
резервирование критичных узлов:
\begin{itemize}[noitemsep]
    \item \textbf{PostgreSQL} --- репликация master\textendash replica
          (streaming replication). При отказе мастера происходит
          автоматическое переключение на реплику с помощью
          \textit{pg\_auto\_failover}. RTO --- не более 30 секунд,
          RPO --- не более 1 секунды.
    \item \textbf{Redis} --- режим Sentinel или Redis Cluster
          (на горизонте развития). На момент защиты используется
          одиночный экземпляр с регулярными снимками AOF.
    \item \textbf{Go API} --- горизонтальное масштабирование
          в нескольких инстансах за общим reverse proxy
          (Caddy с балансировкой). Stateless-приложение,
          состояние сессий в Redis.
    \item \textbf{Резервное копирование} БД --- ежедневные
          логические дампы (\textit{pg\_dump}), хранение
          7~суточных и 4~еженедельных копий.
\end{itemize}

При параллельном включении двух экземпляров узла $N_i$
коэффициент готовности группы:
\[
    k_{\text{г}}^{\,||} = 1 - (1 - k_{\text{г}\, i})^2.
\]
Применение указанной формулы к узлам $N_1$ и $N_2$ повышает
системный коэффициент готовности до $k_{\text{г}}^{\text{сист}\,*}
\geq 0{,}9999$, что соответствует целевому SLA.

\subsection{Идемпотентность операций и согласованность данных}

Для предотвращения двойных записей при ретраях, спровоцированных
сетевыми обрывами, на сервере введены меры:
\begin{itemize}[noitemsep]
    \item уникальное ограничение
          \texttt{UNIQUE (room\_id, question\_id, participant\_id)}
          на таблице \texttt{answers} --- студент не может ответить
          на один и тот же вопрос дважды;
    \item обновление состояния SM\textendash 2 в одной транзакции
          с записью ответа (\texttt{BEGIN; INSERT INTO answers; UPSERT
          INTO sm2\_states; COMMIT}).
    \item отправка ответа клиентом сопровождается заголовком
          \texttt{Idempotency-Key} (UUID запроса); сервер кеширует
          ответ на 5 минут и возвращает то же значение при повторе.
\end{itemize}

\subsection{Безопасность}

\textbf{Канал передачи.} Любой обмен с клиентами шифруется
TLS 1.3 (Caddy с автоматическим выпуском сертификата Let's Encrypt
или ACME-провайдера в РФ-инфраструктуре).

\textbf{Аутентификация.} Используется JWT-токен с временем жизни
15 минут (access) и 7 суток (refresh). Refresh-токены ротируются
при каждом обновлении и инвалидируются на сервере (хранятся в Redis).
Пароли хешируются bcrypt с cost = 12.

\textbf{Авторизация.} Каждый защищённый эндпоинт проверяет
не только валидность JWT, но и принадлежность ресурса
(\textit{tenant isolation}): преподаватель имеет доступ только
к своим курсам, студент --- только к комнатам, участником
которых он является. Контроль продублирован уникальными
ограничениями БД.

\textbf{Защита от типовых уязвимостей} OWASP Top 10:
\begin{itemize}[noitemsep]
    \item SQL-инъекции --- параметризованные запросы (\texttt{pgx});
    \item XSS --- экранирование вывода во Flutter
          (по умолчанию все строки в Text-виджете обрабатываются
          как plain text);
    \item CSRF --- SameSite=Strict cookies для refresh-токена,
          Bearer-токен для остальных запросов;
    \item IDOR --- проверка принадлежности ресурса в каждом обработчике;
    \item Mass assignment --- DTO с явным списком полей
          вместо прямой десериализации в модель.
\end{itemize}

\textbf{Rate limiting.} На уровне Caddy установлен лимит
60 запросов в минуту на IP для эндпоинтов аутентификации
и 600 запросов в минуту для остальных, что снижает поверхность
атаки на брутфорс.

\subsection{Нагрузочное тестирование}

Нагрузочное тестирование проводилось инструментом \textbf{k6}
на отдельном экземпляре сервиса (2 vCPU, 4 ГБ RAM, Selectel)
по сценариям:
\begin{enumerate}[label=\arabic*., noitemsep]
    \item подключение $N$ студентов к одной комнате
          через WebSocket;
    \item массовая подача ответа в момент истечения таймера
          (worst case --- все студенты отправляют ответ
          одновременно);
    \item смена вопроса преподавателем с трансляцией события
          $N$ участникам.
\end{enumerate}

% TODO: вставить таблицу результатов и графики после проведения
% реального нагрузочного теста на стенде. Целевые показатели:

Целевые показатели по результатам тестирования:
\begin{itemize}[noitemsep]
    \item p95 задержки доставки события <<новый вопрос>> ---
          не более 1 секунды при $N = 100$;
    \item p95 времени ответа на API \texttt{POST /answers} ---
          не более 500 мс при $N = 500$;
    \item доля ошибок (HTTP 5xx) --- не более 0{,}1\%
          в условиях штатной нагрузки.
\end{itemize}

\subsection{Мониторинг и алертинг}

Метрики собираются Prometheus и визуализируются в Grafana.
Источники метрик: \textbf{встроенный} HTTP-эндпоинт \texttt{/metrics}
(стандартные метрики Go runtime, gin\_request\_duration\_seconds);
\textbf{postgres\_exporter} (соединения, blocked queries,
replication lag); \textbf{redis\_exporter}; \textbf{node\_exporter}
(CPU, RAM, диск, сеть VM).

Сформированы алерты в Alertmanager:
\begin{itemize}[noitemsep]
    \item рост 5xx-ошибок API выше 1\% за 5 минут;
    \item превышение p95 latency 1 секунды на ключевых эндпоинтах;
    \item replication lag PostgreSQL свыше 5 секунд;
    \item память Redis свыше 80\% от лимита;
    \item недоступность \texttt{/healthz} 3 проверки подряд.
\end{itemize}

\textbf{Выводы по главе 4.}
В главе выполнен расчёт показателей надёжности архитектуры
сервиса EduQuiz: вероятности безотказной работы, средней наработки
на отказ, коэффициента готовности; обосновано резервирование
PostgreSQL master\textendash replica, обеспечивающее системный
коэффициент готовности $k_{\text{г}}^{\,*} \geq 0{,}9999$.
Реализованы меры информационной безопасности (TLS 1.3, JWT,
защита от OWASP Top 10), идемпотентность операций записи,
непрерывный мониторинг и алертинг. Полученные результаты
демонстрируют практическое применение положений курса
<<Надёжность>>, преподаваемого научным руководителем работы,
к собственной архитектуре сервиса.

\newpage
